\chapter{Spine}

\label{ch:ch40_spine}

Back pain is one of the most common reasons for medical imaging. Your doctor may order an X-ray, CT scan, or MRI of the spine depending on the suspected cause. An \textbf{X-ray} is quick and shows bones well---it can reveal fractures, alignment problems, and wear-and-tear changes. A \textbf{CT scan} gives more detailed bone images and is especially useful for planning surgery or detecting subtle fractures. An \textbf{MRI} uses no X-rays and is the best test for looking at the spinal cord, nerve roots, discs, and soft tissues. MRI can show disc herniations pressing on nerves, infection, tumors, and inflammation of the spinal cord.

In emergencies---such as sudden weakness in the legs, loss of bladder control, or severe trauma---urgent imaging is essential to determine whether the spinal cord or nerves are being compressed and whether surgery is needed.


\section{Modality Selection}
Radiography may assess alignment and degenerative change; CT is useful for bone and acute trauma; MRI is preferred for cord, nerve roots, marrow, infection, and many tumors. Choose the smallest appropriate coverage that answers the clinical question.

\begin{table}[h]
\centering
\caption{CT vs MRI for spine imaging}
\begin{tabular}{p{3.5cm} p{5.5cm} p{5.5cm}}
\hline
\textbf{Feature} & \textbf{CT} & \textbf{MRI} \\
\hline
Disc herniation & Limited (no cord detail) & Excellent \\
Fracture detection & Excellent & Excellent \\
Infection / osteomyelitis & Limited sensitivity & Excellent (edema, abscess, discitis) \\
Post-operative assessment & Good for hardware & Better for recurrent disc, infection \\
Radiation & Yes & None \\
Speed & Fast (seconds to minutes) & 20--45 minutes \\
Spinal cord pathology & Limited & Excellent (myelopathy, MS, tumor) \\
Spinal stenosis & Good (bony canal) & Excellent (soft tissue and cord) \\
\hline
\end{tabular}
\end{table}

\begin{pearlbox}[MRI for neurological deficit]
Any patient with progressive neurological symptoms (weakness, numbness, bladder/bowel dysfunction) requires MRI of the affected spinal region as an urgent investigation. Do not delay MRI for CT or X-ray when cord compression is suspected.
\end{pearlbox}


\section{Degenerative Disease}
Degenerative spine disease is nearly universal with ageing. Key findings include:

\begin{itemize}
  \item \textbf{Disc desiccation}: Loss of T2 signal in the nucleus pulposus indicates dehydration and early degeneration.
  \item \textbf{Disc height loss}: Seen on X-ray and CT; correlates with chronic pain.
  \item \textbf{Osteophytes}: Bony spurs at vertebral endplates and facet joints.
  \item \textbf{Facet joint hypertrophy}: Enlarged facets with joint space narrowing and subchondral sclerosis.
  \item \textbf{Endplate changes (Modic types)}:
    \begin{itemize}
      \item Type 1: T1 hypointense, T2 hyperintense (active inflammation)
      \item Type 2: T1 hyperintense, T2 iso/hyperintense (fatty replacement)
      \item Type 3: T1 and T2 hypointense (sclerosis)
    \end{itemize}
\end{itemize}


\section{Disc Herniation}
Disc herniation is classified by morphology:

\begin{table}[h]
\centering
\caption{Disc herniation types}
\begin{tabular}{p{3cm} p{4.5cm} p{4.5cm}}
\hline
\textbf{Type} & \textbf{Definition} & \textbf{Clinical significance} \\
\hline
Broad-based & Herniation involves $> 25\%$ of disc circumference & May cause central canal stenosis \\
Extrusion & Disc material breaches annulus but remains continuous with parent disc & More likely to cause radiculopathy \\
Protrusion & Broad-based; disc material remains within annulus & May cause symptoms if compressing nerve root \\
Sequestration & Free fragment detached from parent disc & May resorb spontaneously; often causes acute radiculopathy \\
\hline
\end{tabular}
\end{table}

Key locations and affected structures:
\begin{itemize}
  \item \textbf{Central herniation}: Compresses the cord (cervical) or cauda equina (lumbar)
  \item \textbf{Paracentral herniation}: Compresses traversing nerve root (most common site for lumbar herniation)
  \item \textbf{Foraminal herniation}: Compresses the exiting nerve root
  \item \textbf{Far lateral (extraforaminal)}: Compresses the exiting nerve root outside the foramen
\end{itemize}


\section{Spinal Stenosis}
Central canal stenosis is graded as:
\begin{itemize}
  \item \textbf{Mild}: Thecal sac compression without CSF effacement
  \item \textbf{Moderate}: CSF effacement but cord/rootlets visible
  \item \textbf{Severe}: Complete obliteration of CSF space with cord/rootlet compression and signal change
\end{itemize}

\textbf{Neuroforaminal stenosis} is assessed by narrowing of the neural foramen due to osteophyte, disc bulge, or facet hypertrophy, causing radiculopathy.


\section{Spine Fractures}
Fracture classification guides management:

\begin{table}[h]
\centering
\caption{Vertebral fracture types}
\begin{tabular}{p{2.5cm} p{3cm} p{3.5cm} p{3.5cm}}
\hline
\textbf{Type} & \textbf{Mechanism} & \textbf{CT features} & \textbf{Stability} \\
\hline
Burst & High-energy axial load & Posterior cortex disrupted; retropulsed fragment into canal & Unstable \\
Chance (flexion--distraction) & Seat-belt injury & Horizontal fracture through body, pedicles, and posterior elements & Unstable \\
Compression & Axial load, flexion & Anterior wedge; posterior cortex intact & Usually stable \\
Fracture--dislocation & Multi-column failure & Vertebral translation, complete disruption & Highly unstable \\
\hline
\end{tabular}
\end{table}

\begin{pearlbox}[Three-column spine]
Denis classified the spine into anterior (anterior longitudinal ligament, anterior annulus, anterior two-thirds of vertebral body), middle (posterior annulus, posterior vertebral body, posterior longitudinal ligament), and posterior (pedicles, laminae, facets, spinous process) columns. Involvement of two or more columns indicates instability.
\end{pearlbox}


\section{Cord Compression}
\textbf{Urgent finding}. Cord compression from disc herniation, tumor, epidural abscess, or fracture requires immediate neurosurgical consultation. MRI is the investigation of choice.

Signs of cord compression on MRI:
\begin{itemize}
  \item Epidural mass (tumor, abscess, disc fragment) abutting the cord
  \item Cord signal change: T2 hyperintensity indicating edema or myelomalacia
  \item Complete obliteration of the CSF space around the cord
\end{itemize}

\begin{pearlbox}[Cauda equina syndrome]
Bilateral leg weakness, saddle anaesthesia, and urinary retention constitute cauda equina syndrome---a neurosurgical emergency. Urgent MRI of the lumbar spine is required. Do not wait for a formal radiology report if clinical suspicion is high; escalate immediately.
\end{pearlbox}


\section{Cord Pathology}

\subsection{Multiple Sclerosis (MS)}
Spinal cord MS plaques appear as ovoid T2-hyperintense lesions, typically short-segment ($< 2$ vertebral bodies), peripherally located, and oriented along the long axis of the cord. Active plaques may show gadolinium enhancement. Lesions at multiple levels are characteristic.

\subsection{Transverse Myelitis}
Longitudinally extensive transverse myelitis (LETM) is defined as a T2-hyperintense cord lesion extending over $\geq 3$ vertebral segments. It is associated with neuromyelitis optica spectrum disorder (NMOSD) and requires different treatment from MS.


\section{Spinal Tumors}
\begin{table}[h]
\centering
\caption{Spinal tumor classification}
\begin{tabular}{p{3cm} p{3.5cm} p{3.5cm} p{3cm}}
\hline
\textbf{Category} & \textbf{Common types} & \textbf{MRI features} & \textbf{Age group} \\
\hline
Benign & Osteoid osteoma, osteoblastoma & Focal nidus, reactive sclerosis & Young adults \\
Extradural & Metastases, lymphoma & Vertebral body replacement, epidural soft tissue & Adults (metastases) \\
Intradural--extramedullary & Meningioma, schwannoma & Well-defined, enhancing, dural tail (meningioma) & Adults \\
Intradural--intramedullary & Ependymoma, astrocytoma & Cord expansion, T2 hyperintensity, enhancement & Variable \\
\hline
\end{tabular}
\end{table}


\section{Review}
Assess alignment, vertebral body height, discs, facets, canal, foramina, cord, marrow, paraspinal soft tissues, and relevant vessels. Report level, side, severity, and neurologic or compressive consequences.

\section{Urgent Findings}
Cord compression, cauda equina compression, epidural abscess, unstable fracture, and destructive lesion require prompt communication. Imaging does not replace neurologic examination or surgical assessment.

\clearpage
